\newpage
\section*{APÊNDICE A --- Definições das medidas e convenções terminológicas}
\addcontentsline{toc}{section}{APÊNDICE A}

\subsection*{Convenções terminológicas}

Resumo da convenção fixada na seção 3.1: em português, \emph{processo} designa ao mesmo tempo o
caso administrativo e o objeto da disciplina de \emph{mineração de processos}. Para evitar a
ambiguidade, este trabalho fixa: \textbf{processo administrativo} é o caso individual;
\textbf{tipo de processo} é o agrupamento emergente; \textbf{mineração de processos} é a
disciplina. \emph{Trace} é traduzido por \textbf{traço}, e \emph{event log} por \textbf{registros
de atividades}. Dos três alfabetos de traço, \textbf{ação por unidade} é o alfabeto não projetado, em que
cada estado é o par ação@unidade e que carrega ambos os eixos. Quanto ao segundo alfabeto,
distingue-se \textbf{roteamento} --- o alfabeto de traço, isto é, a representação modelada --- de
\textbf{tramitação}, o movimento administrativo propriamente dito a que a norma se refere.
\textbf{Procedimento}, o terceiro, conserva apenas a sequência de ações, abstraída do organograma.
Por fim, \textbf{ação} é o tipo de andamento registrado pelo sistema, e \textbf{atividade} é a
ocorrência individual de uma ação em um processo administrativo.

\subsection*{Glossário das medidas}

%% GENERATED by `pipeline run results` - do not edit by hand.
\begin{table}[htbp]
\centering
\caption{Medidas utilizadas neste trabalho.}\label{tab:glossary}
\footnotesize
\resizebox{\textwidth}{!}{%
\begin{tabular}{llcp{0.52\textwidth}}
\toprule
medida & categoria & melhor & definição \\
\midrule
\texttt{top-1} & a jusante & $\uparrow$ & Parcela dos passos retidos em que a primeira predição do modelo é o verdadeiro estado seguinte. \\
\texttt{top-3} & a jusante & $\uparrow$ & Parcela dos passos retidos em que o verdadeiro estado seguinte está entre as três primeiras predições do modelo. \\
\texttt{top-1 (common set)} & a jusante & $\uparrow$ & top-1 recalculado apenas sobre os processos que todo agrupamento avalia, removendo o confundimento de cobertura. Nas tabelas: ``t-1 comum''. \\
\texttt{successor baseline} & linha de base & $\uparrow$ & Consulta ao sucessor mais frequente dentro do grupo. Piso para 'o modelo aprendeu alguma coisa'. \\
\texttt{global model} & controle & $\uparrow$ & Uma única cadeia de Markov de terceira ordem treinada sobre o corpus agregado. Responde 'um único modelo se sairia igualmente bem?'. \\
\texttt{delta-star} & a jusante & $\uparrow$ & O estimando: o escore da partição real menos o escore de seu controle embaralhado de forma equivalente (mesmo número de grupos, mesma distribuição de tamanhos, mesma população, com a pertença permutada). É positivo somente se a partição carrega sinal sobre o processo além da sua forma. Calculado por família de modelo, de modo que a comparação é termo a termo. delta\_star\_adjusted subtrai adicionalmente a linha de base global sob ambas as partições, cancelando o artefato de ponderação que global\_baseline\_residual reporta. É um CRITÉRIO DE ADMISSÃO, não uma ordenação: delta-star cresce com o número de famílias, porque especializar compensa mais quando as famílias são pequenas, e o controle preserva a forma mas não o que a própria granularidade compra. Nunca ordene partições de granularidade diferente por ele -- use selection\_score para isso; comparações em k pareado permanecem válidas. \\
\texttt{global baseline residual} & controle & $\downarrow$ & O modelo global agregado pontuado sobre os mesmos passos retidos sob a partição real e sob a embaralhada. Deve ser \textasciitilde{}0 por construção; um valor diferente de zero significa que a ponderação por n\_traces está produzindo efeito, e o desequilíbrio entre famílias deve ser reportado ao lado de delta-star. \\
\texttt{shuffled partition} & controle & $\uparrow$ & Mesma contagem de grupos, mesmos tamanhos, mesma população, com a pertença permutada. Piso para 'a partição é real?'. Isoladamente é o ganho gratuito de particionar, não um resultado; sua distância até o braço real é o delta-star. \\
\texttt{gain vs global} & a jusante & $\uparrow$ & Acurácia por grupo menos a acurácia da linha de base sobre os mesmos passos retidos. Medido dentro de um único alfabeto contra uma linha de base que enfrenta o mesmo vocabulário, de modo que é a única grandeza comparável entre representações. Ainda contém o ganho que qualquer partição da mesma distribuição de tamanhos obtém de graça, porque ajustar por família reduz o alfabeto de cada modelo — reporte-o ao lado do delta-star, nunca em seu lugar. \\
\texttt{suffix distance} & fidelidade & $\downarrow$ & Damerau-Levenshtein normalizado entre o sufixo que o modelo gera a partir de um meio-prefixo e o sufixo real. A razão de comprimento indica se um modelo vence por parar cedo demais. \\
\texttt{suffix distance, constrained} & fidelidade & $\downarrow$ & A mesma medida sob decodificação restrita, em que o autômato de fluxo filtra todo candidato. Emitida apenas para alfabetos cujos estados carregam uma ação; em um alfabeto de roteamento não há regra de ciclo de vida a impor. \\
\texttt{flow validity} & fidelidade & $\uparrow$ & Parcela dos passos gerados (e dos traços gerados inteiramente legais) que o autômato de ciclo de vida do processo permite. O par \_real é a mesma medida sobre os registros de atividades observados, que é o teto. \\
\texttt{unit\_recall} & cobertura & $\uparrow$ & Fração das unidades visitadas por um processo real que o caminho gerado também visita. Nas tabelas: ``revocação de unidades''. \\
\texttt{unit\_precision} & cobertura & $\uparrow$ & Fração das unidades no caminho gerado que os processos reais de fato visitam. \\
\texttt{visits\_matched\_frac} & cobertura & $\uparrow$ & Fração das visitas de um processo real correspondidas pelo caminho gerado. Capta a cobertura no nível do passo, não apenas quais unidades. \\
\texttt{transition\_support} & cobertura & $\uparrow$ & Número médio de ocorrências de cada transição gerada por membro do grupo. Conta ocorrências, não membros distintos, de modo que pode exceder 1 e não é uma proporção. A variante \_min é o elo mais fraco do caminho. \\
\texttt{transitions\_unseen} & fidelidade & $\downarrow$ & Transições do caminho gerado que não ocorrem em nenhum traço retido amostrado do grupo, contra uma amostra de no máximo trinta. Fabricação, não variação; a medida a superestima para grupos maiores do que essa amostra. \\
\texttt{unseen per visit} & fidelidade & $\downarrow$ & transitions\_unseen normalizado pelo comprimento do caminho, de modo que um caminho curto não pareça fiel por gerar menos. A chave simples é por grupo; a chave \_pooled é a razão, no nível da população, entre médias ponderadas, e as duas não são intercambiáveis. \\
\texttt{within\_25} & fidelidade & $\uparrow$ & Parcela dos traços retidos a até 0,25 de Damerau-Levenshtein normalizado do caminho gerado. \\
\texttt{typical path distance} & fidelidade & $\downarrow$ & Damerau-Levenshtein normalizado médio entre o caminho típico gerado e os traços retidos amostrados. As variantes \_proc e \_itin projetam sobre outro alfabeto e existem apenas para células cujos estados carregam os dois eixos. \\
\texttt{perplexity} & a jusante & $\downarrow$ & Exponencial da média da log-probabilidade negativa que o modelo atribui ao verdadeiro próximo passo: o número efetivo de opções entre as quais ele ainda está decidindo. oov\_rate é a parcela de passos aos quais ele atribuiu probabilidade zero e que foram elevados ao piso. \\
\texttt{edge coverage} & fidelidade & $\uparrow$ & Parcela das transições retidas presentes no grafo de precedência direta não podado do grupo; edge\_usage é a recíproca, a parcela das arestas do grafo que os traços retidos usam. Ambas são estatísticas de célula calculadas com piso 0,0, independentes da poda do próprio modelo dfg. edge\_fidelity carrega as mesmas três medidas contrastadas contra o controle embaralhado (real/shuffled/delta\_star); como descrevem o grupo e não um modelo, assumem um único valor por partição e não podem variar por família de modelo. \\
\texttt{dispersion} & contexto & --- & Distância média entre pares dentro de um grupo. Limita o quão próximo um único caminho pode estar de todos os membros. As variantes \_proc e \_itin medem a mesma dispersão após a projeção sobre outro alfabeto, e existem apenas para células cujos estados carregam os dois eixos. \\
\texttt{medoid} & linha de base & --- & O traço real mais próximo, em média, dos demais. Linha de base de proximidade contra a qual o caminho gerado precisa se justificar. \\
\texttt{DBCV} & validade da partição & $\uparrow$ & Validade de grupo baseada em densidade. Concebida para formas arbitrárias; severa com partições por centroide. \\
\texttt{silhouette} & validade da partição & $\uparrow$ & Coesão vs separação. Premia grupos compactos e convexos — a forma que o k-means otimiza diretamente. \\
\texttt{Davies-Bouldin} & validade da partição & $\downarrow$ & Razão média de pior caso entre a dispersão intragrupo e a separação entre grupos, no espaço cosseno UMAP-50. \\
\texttt{Calinski-Harabasz} & validade da partição & $\uparrow$ & Razão entre a dispersão entre grupos e a dispersão intragrupo, no espaço cosseno UMAP-50. \\
\texttt{cluster separation} & validade da partição & $\uparrow$ & Distância ao centroide mais próximo sobre o raio médio do próprio grupo, com média entre os grupos. Abaixo de 1,0 os grupos se tocam; abaixo de 0,7 eles se sobrepõem. \\
\texttt{Jaccard stability} & validade da partição & $\uparrow$ & Recuperação de grupos por bootstrap em 25 reajustes de subamostra. Limite inferior: o UMAP não é reajustado. \\
\texttt{\% evaluated} & cobertura & $\uparrow$ & Parcela do corpus dentro de um grupo que sobrevive ao piso min\_traces. Capta AMBOS os descartes. \\
\texttt{noise ratio} & cobertura & $\downarrow$ & Parcela que um método por densidade deixa sem atribuição. noise\_ratio é o que resta APÓS a atribuição padrão ao grupo mais próximo, de modo que é zero em toda parte e não pode ser correlacionado; noise\_ratio\_refused (armazenado por agrupamento como noise\_ratio\_before\_assignment) é o que a densidade de fato recusou, e é a variável que toda alegação de abstenção utiliza. Zero para k-means e Leiden por construção. Nas tabelas: ``ruído''. \\
\texttt{ARI / NMI vs tipo} & descritivo & --- & Concordância com a taxonomia declarada. Apenas reportada; não utilizada como critério de seleção. \\
\texttt{coverage x fidelity} & composto & $\uparrow$ & unit\_recall x top-1 / (1 + transitions\_unseen). Composto explícito usado como variável de resultado da correlação validade-vs-utilidade e como chave de ordenação das tabelas de combinações, de modo que um único número contrapõe cobertura e fabricação. Seus componentes são sempre reportados separadamente; ele não tem existência independente. \\
\texttt{selection score} & composto & $\uparrow$ & Combinação escalonada por mín--máx de acurácia no conjunto comum, cobertura e silhouette usada para eleger a partição entregável; reportada sob quatro ponderações, e a eleição precisa sobreviver a todas elas. \\
\texttt{selected-best top-1} & a jusante & $\uparrow$ & Máximo por grupo entre todas as famílias de modelo. Não é comparação equivalente contra a linha de base global de família única; os ganhos com família fixa são reportados ao lado. Nas tabelas: ``melhor sel.''. \\
\texttt{non-inferiority} & controle & $\uparrow$ & Bootstrap pareado, no nível do processo, do braço anonimizado menos o braço bruto, contra a margem pré-registrada. O pareamento é por protocolo dentro de um mesmo algoritmo; a divisão é atribuída para todo o corpus, de modo que ambos os braços pontuam os mesmos processos sempre que ambos os agruparam. \\
\texttt{weighting} & contexto & --- & Todo agregado é reportado sob as duas ponderações: macro/unweighted dá um voto a cada grupo; weighted/membership\_weighted pondera pelo tamanho do grupo. Um ganho que sobrevive apenas a uma delas é uma propriedade da ponderação, não do método. \\
\texttt{constrained decoding cost} & controle & --- & Com que frequência o autômato de fluxo de fato altera os candidatos ranqueados, a primeira predição, ou não deixa nenhum candidato legal. \\
\bottomrule
\end{tabular}
}
\end{table}


\newpage
\section*{APÊNDICE B --- Tabelas complementares}
\addcontentsline{toc}{section}{APÊNDICE B}

As tabelas a seguir complementam o audit trail do capítulo 4 sem repetir os mesmos artefatos já
apresentados no corpo do texto. A grade completa --- as catorze partições avaliadas por família de
modelo e alfabeto, 84 linhas por alfabeto --- sustenta a contagem de células declarada no capítulo 4.

%% GENERATED by `pipeline run results` - do not edit by hand.
{\footnotesize
\begin{longtable}{llrrrrr}
\caption{Grade completa de resultados, alfabeto de ação por unidade.}\label{tab:combos-unit}\\
\toprule
agrupamento & família & grupos & top-1 & revocação & precisão & inéditas \\
\midrule
\endfirsthead
\toprule
agrupamento & família & grupos & top-1 & revocação & precisão & inéditas \\
\midrule
\endhead
\bottomrule
\endfoot
exp-anon-hcoarse/hdbscan & transformer & 33 & 0,725 & 0,685 & 0,711 & 2,808 \\
exp-anon-hcoarse/hdbscan & lstm & 33 & 0,715 & 0,681 & 0,691 & 2,876 \\
exp-anon-hcoarse/hdbscan & gru & 33 & 0,726 & 0,669 & 0,699 & 3,372 \\
exp-anon-hcoarse/hdbscan & markov & 33 & 0,720 & 0,620 & 0,710 & 3,604 \\
exp-anon-hcoarse/hdbscan & dfg & 33 & 0,571 & 0,603 & 0,710 & 5,043 \\
exp-anon-hcoarse/hdbscan & gnn & 33 & 0,619 & 0,634 & 0,700 & 5,930 \\
exp-anon-kelbow/kmeans & transformer & 55 & 0,711 & 0,727 & 0,757 & 2,620 \\
exp-anon-kelbow/kmeans & gru & 55 & 0,716 & 0,716 & 0,742 & 2,922 \\
exp-anon-kelbow/kmeans & lstm & 55 & 0,703 & 0,730 & 0,735 & 3,229 \\
exp-anon-kelbow/kmeans & dfg & 55 & 0,588 & 0,655 & 0,760 & 2,236 \\
exp-anon-kelbow/kmeans & markov & 55 & 0,712 & 0,678 & 0,772 & 3,369 \\
exp-anon-kelbow/kmeans & gnn & 55 & 0,624 & 0,662 & 0,777 & 4,193 \\
exp-anon-ksil/kmeans & transformer & 49 & 0,709 & 0,735 & 0,745 & 2,482 \\
exp-anon-ksil/kmeans & gru & 49 & 0,712 & 0,706 & 0,750 & 2,446 \\
exp-anon-ksil/kmeans & lstm & 49 & 0,702 & 0,716 & 0,750 & 3,512 \\
exp-anon-ksil/kmeans & markov & 49 & 0,709 & 0,667 & 0,748 & 3,362 \\
exp-anon-ksil/kmeans & dfg & 49 & 0,581 & 0,677 & 0,761 & 2,646 \\
exp-anon-ksil/kmeans & gnn & 49 & 0,619 & 0,670 & 0,776 & 4,680 \\
exp-anon-lres/leiden & transformer & 78 & 0,707 & 0,763 & 0,807 & 2,540 \\
exp-anon-lres/leiden & dfg & 78 & 0,610 & 0,712 & 0,810 & 2,048 \\
exp-anon-lres/leiden & gru & 78 & 0,711 & 0,743 & 0,801 & 2,882 \\
exp-anon-lres/leiden & lstm & 78 & 0,700 & 0,744 & 0,801 & 3,870 \\
exp-anon-lres/leiden & gnn & 78 & 0,634 & 0,701 & 0,817 & 3,670 \\
exp-anon-lres/leiden & markov & 78 & 0,708 & 0,726 & 0,792 & 4,820 \\
exp-anon/hdbscan & gru & 54 & 0,716 & 0,725 & 0,771 & 1,816 \\
exp-anon/hdbscan & transformer & 54 & 0,715 & 0,735 & 0,762 & 2,107 \\
exp-anon/hdbscan & lstm & 54 & 0,704 & 0,721 & 0,750 & 2,946 \\
exp-anon/hdbscan & dfg & 54 & 0,585 & 0,666 & 0,784 & 2,496 \\
exp-anon/hdbscan & markov & 54 & 0,714 & 0,668 & 0,785 & 3,295 \\
exp-anon/hdbscan & gnn & 54 & 0,624 & 0,661 & 0,795 & 3,930 \\
exp-anon/kmeans & gru & 52 & 0,716 & 0,722 & 0,762 & 2,189 \\
exp-anon/kmeans & transformer & 52 & 0,711 & 0,722 & 0,749 & 2,263 \\
exp-anon/kmeans & lstm & 52 & 0,706 & 0,729 & 0,744 & 3,255 \\
exp-anon/kmeans & dfg & 52 & 0,586 & 0,678 & 0,773 & 2,617 \\
exp-anon/kmeans & markov & 52 & 0,714 & 0,664 & 0,762 & 3,450 \\
exp-anon/kmeans & gnn & 52 & 0,624 & 0,669 & 0,770 & 4,421 \\
exp-anon/leiden & transformer & 55 & 0,718 & 0,724 & 0,757 & 2,279 \\
exp-anon/leiden & gru & 55 & 0,719 & 0,728 & 0,769 & 2,703 \\
exp-anon/leiden & dfg & 55 & 0,593 & 0,683 & 0,786 & 2,041 \\
exp-anon/leiden & lstm & 55 & 0,710 & 0,705 & 0,773 & 2,826 \\
exp-anon/leiden & markov & 55 & 0,716 & 0,680 & 0,771 & 4,316 \\
exp-anon/leiden & gnn & 55 & 0,630 & 0,656 & 0,774 & 4,392 \\
exp-raw-hcoarse/hdbscan & transformer & 36 & 0,718 & 0,676 & 0,712 & 1,791 \\
exp-raw-hcoarse/hdbscan & lstm & 36 & 0,711 & 0,689 & 0,722 & 2,274 \\
exp-raw-hcoarse/hdbscan & dfg & 36 & 0,572 & 0,621 & 0,718 & 2,442 \\
exp-raw-hcoarse/hdbscan & gru & 36 & 0,719 & 0,677 & 0,708 & 3,823 \\
exp-raw-hcoarse/hdbscan & markov & 36 & 0,714 & 0,625 & 0,719 & 3,601 \\
exp-raw-hcoarse/hdbscan & gnn & 36 & 0,616 & 0,638 & 0,714 & 4,054 \\
exp-raw-kelbow/kmeans & transformer & 57 & 0,708 & 0,734 & 0,798 & 1,748 \\
exp-raw-kelbow/kmeans & lstm & 57 & 0,698 & 0,734 & 0,800 & 2,366 \\
exp-raw-kelbow/kmeans & dfg & 57 & 0,592 & 0,699 & 0,809 & 2,376 \\
exp-raw-kelbow/kmeans & markov & 57 & 0,708 & 0,706 & 0,793 & 3,556 \\
exp-raw-kelbow/kmeans & gru & 57 & 0,710 & 0,741 & 0,773 & 4,509 \\
exp-raw-kelbow/kmeans & gnn & 57 & 0,624 & 0,703 & 0,804 & 5,277 \\
exp-raw-ksil/kmeans & lstm & 50 & 0,700 & 0,730 & 0,761 & 2,361 \\
exp-raw-ksil/kmeans & transformer & 50 & 0,706 & 0,729 & 0,769 & 2,508 \\
exp-raw-ksil/kmeans & gru & 50 & 0,710 & 0,712 & 0,753 & 3,173 \\
exp-raw-ksil/kmeans & markov & 50 & 0,710 & 0,665 & 0,758 & 3,535 \\
exp-raw-ksil/kmeans & dfg & 50 & 0,589 & 0,686 & 0,769 & 3,949 \\
exp-raw-ksil/kmeans & gnn & 50 & 0,622 & 0,695 & 0,778 & 4,928 \\
exp-raw-lres/leiden & transformer & 78 & 0,709 & 0,743 & 0,788 & 2,224 \\
exp-raw-lres/leiden & lstm & 78 & 0,702 & 0,749 & 0,796 & 3,365 \\
exp-raw-lres/leiden & dfg & 78 & 0,609 & 0,711 & 0,802 & 2,709 \\
exp-raw-lres/leiden & gru & 78 & 0,712 & 0,742 & 0,791 & 3,950 \\
exp-raw-lres/leiden & markov & 78 & 0,711 & 0,706 & 0,780 & 3,949 \\
exp-raw-lres/leiden & gnn & 78 & 0,633 & 0,707 & 0,791 & 3,751 \\
exp-raw/hdbscan & transformer & 56 & 0,712 & 0,717 & 0,773 & 1,986 \\
exp-raw/hdbscan & lstm & 56 & 0,707 & 0,709 & 0,764 & 2,764 \\
exp-raw/hdbscan & gru & 56 & 0,715 & 0,715 & 0,772 & 3,184 \\
exp-raw/hdbscan & dfg & 56 & 0,587 & 0,669 & 0,779 & 2,418 \\
exp-raw/hdbscan & markov & 56 & 0,712 & 0,665 & 0,769 & 3,902 \\
exp-raw/hdbscan & gnn & 56 & 0,627 & 0,680 & 0,768 & 4,220 \\
exp-raw/kmeans & transformer & 52 & 0,709 & 0,745 & 0,753 & 1,694 \\
exp-raw/kmeans & lstm & 52 & 0,700 & 0,745 & 0,757 & 2,612 \\
exp-raw/kmeans & gru & 52 & 0,711 & 0,709 & 0,751 & 2,856 \\
exp-raw/kmeans & gnn & 52 & 0,623 & 0,710 & 0,764 & 3,831 \\
exp-raw/kmeans & dfg & 52 & 0,588 & 0,686 & 0,770 & 3,435 \\
exp-raw/kmeans & markov & 52 & 0,709 & 0,666 & 0,748 & 4,217 \\
exp-raw/leiden & transformer & 59 & 0,717 & 0,713 & 0,763 & 1,530 \\
exp-raw/leiden & gru & 59 & 0,720 & 0,711 & 0,760 & 2,648 \\
exp-raw/leiden & lstm & 59 & 0,711 & 0,713 & 0,776 & 3,137 \\
exp-raw/leiden & dfg & 59 & 0,595 & 0,694 & 0,773 & 2,698 \\
exp-raw/leiden & markov & 59 & 0,716 & 0,655 & 0,756 & 4,220 \\
exp-raw/leiden & gnn & 59 & 0,631 & 0,667 & 0,777 & 3,763 \\
\end{longtable}
}

{\footnotesize
\begin{longtable}{llrrrrr}
\caption{Grade completa de resultados, alfabeto de roteamento.}\label{tab:combos-itinerary}\\
\toprule
agrupamento & família & grupos & top-1 & revocação & precisão & inéditas \\
\midrule
\endfirsthead
\toprule
agrupamento & família & grupos & top-1 & revocação & precisão & inéditas \\
\midrule
\endhead
\bottomrule
\endfoot
exp-anon-hcoarse/hdbscan & gru & 28 & 0,706 & 0,665 & 0,750 & 0,178 \\
exp-anon-hcoarse/hdbscan & markov & 28 & 0,698 & 0,646 & 0,739 & 0,159 \\
exp-anon-hcoarse/hdbscan & lstm & 28 & 0,687 & 0,666 & 0,743 & 0,279 \\
exp-anon-hcoarse/hdbscan & transformer & 28 & 0,713 & 0,670 & 0,733 & 1,018 \\
exp-anon-hcoarse/hdbscan & gnn & 28 & 0,466 & 0,650 & 0,702 & 0,450 \\
exp-anon-hcoarse/hdbscan & dfg & 28 & 0,439 & 0,634 & 0,722 & 0,541 \\
exp-anon-kelbow/kmeans & markov & 43 & 0,695 & 0,684 & 0,780 & 0,427 \\
exp-anon-kelbow/kmeans & transformer & 43 & 0,707 & 0,712 & 0,790 & 0,581 \\
exp-anon-kelbow/kmeans & gru & 43 & 0,700 & 0,706 & 0,788 & 0,570 \\
exp-anon-kelbow/kmeans & lstm & 43 & 0,685 & 0,710 & 0,780 & 0,583 \\
exp-anon-kelbow/kmeans & dfg & 43 & 0,459 & 0,677 & 0,771 & 0,474 \\
exp-anon-kelbow/kmeans & gnn & 43 & 0,476 & 0,684 & 0,760 & 0,688 \\
exp-anon-ksil/kmeans & markov & 39 & 0,696 & 0,681 & 0,763 & 0,386 \\
exp-anon-ksil/kmeans & gru & 39 & 0,699 & 0,702 & 0,789 & 0,569 \\
exp-anon-ksil/kmeans & transformer & 39 & 0,706 & 0,709 & 0,764 & 0,753 \\
exp-anon-ksil/kmeans & lstm & 39 & 0,687 & 0,707 & 0,771 & 0,719 \\
exp-anon-ksil/kmeans & dfg & 39 & 0,454 & 0,679 & 0,772 & 0,663 \\
exp-anon-ksil/kmeans & gnn & 39 & 0,473 & 0,694 & 0,753 & 0,783 \\
exp-anon-lres/leiden & transformer & 57 & 0,718 & 0,766 & 0,808 & 0,406 \\
exp-anon-lres/leiden & gru & 57 & 0,716 & 0,774 & 0,824 & 0,522 \\
exp-anon-lres/leiden & markov & 57 & 0,707 & 0,737 & 0,814 & 0,509 \\
exp-anon-lres/leiden & lstm & 57 & 0,702 & 0,761 & 0,819 & 0,764 \\
exp-anon-lres/leiden & dfg & 57 & 0,483 & 0,738 & 0,808 & 0,417 \\
exp-anon-lres/leiden & gnn & 57 & 0,497 & 0,731 & 0,803 & 0,796 \\
exp-anon/hdbscan & gru & 41 & 0,699 & 0,707 & 0,801 & 0,318 \\
exp-anon/hdbscan & markov & 41 & 0,694 & 0,686 & 0,796 & 0,425 \\
exp-anon/hdbscan & transformer & 41 & 0,711 & 0,720 & 0,781 & 0,665 \\
exp-anon/hdbscan & lstm & 41 & 0,683 & 0,703 & 0,786 & 0,819 \\
exp-anon/hdbscan & gnn & 41 & 0,476 & 0,703 & 0,770 & 0,585 \\
exp-anon/hdbscan & dfg & 41 & 0,458 & 0,690 & 0,775 & 0,812 \\
exp-anon/kmeans & gru & 42 & 0,702 & 0,712 & 0,773 & 0,344 \\
exp-anon/kmeans & transformer & 42 & 0,704 & 0,721 & 0,761 & 0,439 \\
exp-anon/kmeans & markov & 42 & 0,695 & 0,679 & 0,773 & 0,420 \\
exp-anon/kmeans & lstm & 42 & 0,686 & 0,715 & 0,772 & 0,706 \\
exp-anon/kmeans & dfg & 42 & 0,458 & 0,683 & 0,770 & 0,415 \\
exp-anon/kmeans & gnn & 42 & 0,475 & 0,700 & 0,757 & 0,647 \\
exp-anon/leiden & gru & 43 & 0,709 & 0,735 & 0,795 & 0,289 \\
exp-anon/leiden & markov & 43 & 0,704 & 0,708 & 0,799 & 0,322 \\
exp-anon/leiden & transformer & 43 & 0,716 & 0,731 & 0,790 & 0,443 \\
exp-anon/leiden & lstm & 43 & 0,690 & 0,727 & 0,784 & 0,893 \\
exp-anon/leiden & dfg & 43 & 0,463 & 0,702 & 0,777 & 0,513 \\
exp-anon/leiden & gnn & 43 & 0,483 & 0,707 & 0,765 & 0,714 \\
exp-raw-hcoarse/hdbscan & markov & 30 & 0,697 & 0,674 & 0,761 & 0,264 \\
exp-raw-hcoarse/hdbscan & transformer & 30 & 0,713 & 0,651 & 0,741 & 0,297 \\
exp-raw-hcoarse/hdbscan & lstm & 30 & 0,694 & 0,665 & 0,752 & 0,295 \\
exp-raw-hcoarse/hdbscan & gru & 30 & 0,706 & 0,682 & 0,752 & 0,379 \\
exp-raw-hcoarse/hdbscan & gnn & 30 & 0,466 & 0,699 & 0,713 & 0,457 \\
exp-raw-hcoarse/hdbscan & dfg & 30 & 0,445 & 0,665 & 0,729 & 0,396 \\
exp-raw-kelbow/kmeans & markov & 43 & 0,700 & 0,716 & 0,823 & 0,124 \\
exp-raw-kelbow/kmeans & transformer & 43 & 0,715 & 0,738 & 0,806 & 0,279 \\
exp-raw-kelbow/kmeans & gru & 43 & 0,708 & 0,748 & 0,824 & 0,344 \\
exp-raw-kelbow/kmeans & lstm & 43 & 0,694 & 0,725 & 0,812 & 0,579 \\
exp-raw-kelbow/kmeans & dfg & 43 & 0,466 & 0,712 & 0,782 & 0,300 \\
exp-raw-kelbow/kmeans & gnn & 43 & 0,481 & 0,730 & 0,781 & 0,746 \\
exp-raw-ksil/kmeans & markov & 39 & 0,696 & 0,683 & 0,791 & 0,262 \\
exp-raw-ksil/kmeans & gru & 39 & 0,705 & 0,725 & 0,791 & 0,491 \\
exp-raw-ksil/kmeans & lstm & 39 & 0,692 & 0,723 & 0,790 & 0,540 \\
exp-raw-ksil/kmeans & transformer & 39 & 0,709 & 0,721 & 0,785 & 0,735 \\
exp-raw-ksil/kmeans & dfg & 39 & 0,458 & 0,718 & 0,773 & 0,672 \\
exp-raw-ksil/kmeans & gnn & 39 & 0,475 & 0,719 & 0,744 & 0,908 \\
exp-raw-lres/leiden & gru & 58 & 0,711 & 0,765 & 0,815 & 0,278 \\
exp-raw-lres/leiden & transformer & 58 & 0,716 & 0,746 & 0,800 & 0,392 \\
exp-raw-lres/leiden & markov & 58 & 0,703 & 0,717 & 0,813 & 0,360 \\
exp-raw-lres/leiden & lstm & 58 & 0,700 & 0,750 & 0,818 & 0,431 \\
exp-raw-lres/leiden & dfg & 58 & 0,484 & 0,732 & 0,789 & 0,337 \\
exp-raw-lres/leiden & gnn & 58 & 0,492 & 0,730 & 0,774 & 0,750 \\
exp-raw/hdbscan & markov & 43 & 0,691 & 0,691 & 0,813 & 0,160 \\
exp-raw/hdbscan & gru & 43 & 0,693 & 0,727 & 0,807 & 0,378 \\
exp-raw/hdbscan & transformer & 43 & 0,700 & 0,708 & 0,792 & 0,429 \\
exp-raw/hdbscan & lstm & 43 & 0,679 & 0,723 & 0,811 & 0,565 \\
exp-raw/hdbscan & dfg & 43 & 0,458 & 0,695 & 0,783 & 0,597 \\
exp-raw/hdbscan & gnn & 43 & 0,472 & 0,712 & 0,763 & 0,999 \\
exp-raw/kmeans & markov & 41 & 0,698 & 0,704 & 0,770 & 0,257 \\
exp-raw/kmeans & gru & 41 & 0,702 & 0,736 & 0,769 & 0,403 \\
exp-raw/kmeans & lstm & 41 & 0,688 & 0,733 & 0,773 & 0,380 \\
exp-raw/kmeans & transformer & 41 & 0,709 & 0,719 & 0,758 & 0,399 \\
exp-raw/kmeans & dfg & 41 & 0,459 & 0,710 & 0,740 & 0,565 \\
exp-raw/kmeans & gnn & 41 & 0,474 & 0,717 & 0,720 & 0,792 \\
exp-raw/leiden & markov & 42 & 0,704 & 0,695 & 0,798 & 0,106 \\
exp-raw/leiden & gru & 42 & 0,710 & 0,729 & 0,789 & 0,189 \\
exp-raw/leiden & lstm & 42 & 0,693 & 0,719 & 0,796 & 0,275 \\
exp-raw/leiden & transformer & 42 & 0,721 & 0,720 & 0,778 & 0,399 \\
exp-raw/leiden & dfg & 42 & 0,466 & 0,696 & 0,773 & 0,228 \\
exp-raw/leiden & gnn & 42 & 0,484 & 0,675 & 0,752 & 0,707 \\
\end{longtable}
}

{\footnotesize
\begin{longtable}{llrrr}
\caption{Grade completa de resultados, alfabeto de ação.}\label{tab:combos-procedure}\\
\toprule
agrupamento & família & grupos & top-1 & inéditas \\
\midrule
\endfirsthead
\toprule
agrupamento & família & grupos & top-1 & inéditas \\
\midrule
\endhead
\bottomrule
\endfoot
exp-anon-hcoarse/hdbscan & markov & 33 & 0,732 & 1,548 \\
exp-anon-hcoarse/hdbscan & dfg & 33 & 0,641 & 0,570 \\
exp-anon-hcoarse/hdbscan & gnn & 33 & 0,642 & 2,151 \\
exp-anon-hcoarse/hdbscan & gru & 33 & 0,756 & 3,644 \\
exp-anon-hcoarse/hdbscan & lstm & 33 & 0,756 & 3,203 \\
exp-anon-hcoarse/hdbscan & transformer & 33 & 0,752 & 2,766 \\
exp-anon-kelbow/kmeans & markov & 55 & 0,730 & 1,674 \\
exp-anon-kelbow/kmeans & dfg & 55 & 0,644 & 0,653 \\
exp-anon-kelbow/kmeans & gnn & 55 & 0,644 & 2,156 \\
exp-anon-kelbow/kmeans & gru & 55 & 0,754 & 3,105 \\
exp-anon-kelbow/kmeans & lstm & 55 & 0,754 & 2,417 \\
exp-anon-kelbow/kmeans & transformer & 55 & 0,750 & 2,134 \\
exp-anon-ksil/kmeans & markov & 49 & 0,731 & 1,421 \\
exp-anon-ksil/kmeans & dfg & 49 & 0,640 & 0,636 \\
exp-anon-ksil/kmeans & gnn & 49 & 0,642 & 1,949 \\
exp-anon-ksil/kmeans & gru & 49 & 0,753 & 2,537 \\
exp-anon-ksil/kmeans & lstm & 49 & 0,751 & 2,153 \\
exp-anon-ksil/kmeans & transformer & 49 & 0,750 & 1,838 \\
exp-anon-lres/leiden & markov & 78 & 0,735 & 1,944 \\
exp-anon-lres/leiden & dfg & 78 & 0,648 & 0,819 \\
exp-anon-lres/leiden & gnn & 78 & 0,649 & 2,586 \\
exp-anon-lres/leiden & gru & 78 & 0,754 & 3,145 \\
exp-anon-lres/leiden & lstm & 78 & 0,754 & 2,693 \\
exp-anon-lres/leiden & transformer & 78 & 0,751 & 2,730 \\
exp-anon/hdbscan & markov & 54 & 0,732 & 2,432 \\
exp-anon/hdbscan & dfg & 54 & 0,643 & 0,712 \\
exp-anon/hdbscan & gnn & 54 & 0,644 & 1,895 \\
exp-anon/hdbscan & gru & 54 & 0,756 & 2,904 \\
exp-anon/hdbscan & lstm & 54 & 0,755 & 3,003 \\
exp-anon/hdbscan & transformer & 54 & 0,752 & 1,727 \\
exp-anon/kmeans & markov & 52 & 0,732 & 1,708 \\
exp-anon/kmeans & dfg & 52 & 0,644 & 0,670 \\
exp-anon/kmeans & gnn & 52 & 0,644 & 2,355 \\
exp-anon/kmeans & gru & 52 & 0,755 & 2,841 \\
exp-anon/kmeans & lstm & 52 & 0,755 & 2,453 \\
exp-anon/kmeans & transformer & 52 & 0,753 & 2,163 \\
exp-anon/leiden & markov & 55 & 0,735 & 1,708 \\
exp-anon/leiden & dfg & 55 & 0,646 & 0,739 \\
exp-anon/leiden & gnn & 55 & 0,647 & 1,691 \\
exp-anon/leiden & gru & 55 & 0,757 & 3,681 \\
exp-anon/leiden & lstm & 55 & 0,759 & 2,528 \\
exp-anon/leiden & transformer & 55 & 0,753 & 2,572 \\
exp-raw-hcoarse/hdbscan & markov & 36 & 0,731 & 2,710 \\
exp-raw-hcoarse/hdbscan & dfg & 36 & 0,640 & 0,613 \\
exp-raw-hcoarse/hdbscan & gnn & 36 & 0,639 & 2,458 \\
exp-raw-hcoarse/hdbscan & gru & 36 & 0,754 & 2,433 \\
exp-raw-hcoarse/hdbscan & lstm & 36 & 0,751 & 5,013 \\
exp-raw-hcoarse/hdbscan & transformer & 36 & 0,749 & 5,787 \\
exp-raw-kelbow/kmeans & markov & 57 & 0,730 & 2,271 \\
exp-raw-kelbow/kmeans & dfg & 57 & 0,642 & 0,892 \\
exp-raw-kelbow/kmeans & gnn & 57 & 0,643 & 2,003 \\
exp-raw-kelbow/kmeans & gru & 57 & 0,753 & 3,211 \\
exp-raw-kelbow/kmeans & lstm & 57 & 0,753 & 2,723 \\
exp-raw-kelbow/kmeans & transformer & 57 & 0,750 & 2,007 \\
exp-raw-ksil/kmeans & markov & 50 & 0,733 & 2,392 \\
exp-raw-ksil/kmeans & dfg & 50 & 0,643 & 0,739 \\
exp-raw-ksil/kmeans & gnn & 50 & 0,643 & 2,133 \\
exp-raw-ksil/kmeans & gru & 50 & 0,754 & 2,498 \\
exp-raw-ksil/kmeans & lstm & 50 & 0,753 & 2,800 \\
exp-raw-ksil/kmeans & transformer & 50 & 0,750 & 1,784 \\
exp-raw-lres/leiden & markov & 78 & 0,734 & 1,846 \\
exp-raw-lres/leiden & dfg & 78 & 0,649 & 0,942 \\
exp-raw-lres/leiden & gnn & 78 & 0,649 & 2,716 \\
exp-raw-lres/leiden & gru & 78 & 0,753 & 2,673 \\
exp-raw-lres/leiden & lstm & 78 & 0,752 & 3,221 \\
exp-raw-lres/leiden & transformer & 78 & 0,751 & 3,269 \\
exp-raw/hdbscan & markov & 56 & 0,736 & 2,342 \\
exp-raw/hdbscan & dfg & 56 & 0,648 & 0,724 \\
exp-raw/hdbscan & gnn & 56 & 0,648 & 2,359 \\
exp-raw/hdbscan & gru & 56 & 0,756 & 3,359 \\
exp-raw/hdbscan & lstm & 56 & 0,754 & 5,039 \\
exp-raw/hdbscan & transformer & 56 & 0,751 & 4,886 \\
exp-raw/kmeans & markov & 52 & 0,731 & 2,196 \\
exp-raw/kmeans & dfg & 52 & 0,642 & 0,862 \\
exp-raw/kmeans & gnn & 52 & 0,641 & 2,285 \\
exp-raw/kmeans & gru & 52 & 0,753 & 3,086 \\
exp-raw/kmeans & lstm & 52 & 0,751 & 2,466 \\
exp-raw/kmeans & transformer & 52 & 0,749 & 2,067 \\
exp-raw/leiden & markov & 59 & 0,735 & 2,236 \\
exp-raw/leiden & dfg & 59 & 0,645 & 0,663 \\
exp-raw/leiden & gnn & 59 & 0,646 & 1,911 \\
exp-raw/leiden & gru & 59 & 0,754 & 2,821 \\
exp-raw/leiden & lstm & 59 & 0,753 & 2,728 \\
exp-raw/leiden & transformer & 59 & 0,751 & 2,245 \\
\end{longtable}
}

